This section introduces the geometric object of the whole series.  An arithmetic
expression space is a space together with an assignment function and a metric on
which the four arithmetic operations act as *motions*: each elementary
operation moves the current arithmetic value, and the effect of that motion on
points is part of the structure.  We give the axioms, the canonical frame they
determine, the basic hyperbolic model, the four elementary motions, the point
paths they generate, and the continuous form of the same motion.

\subsection{Arithmetic motions}

Let $M$ be a smooth surface with a Riemannian metric $g$, and let
\[
  a:M\longrightarrow\mathbb R
\]
be a smooth function, called the \emph{assignment}: $a(p)$ is the arithmetic
value carried at the point $p$.  An arithmetic operation is to be realized as a
motion of $M$ whose effect on $a$ is that operation.  For instance, adding a
constant $s$ should be a map $\mathsf X_s$ with
\[
  a\circ\mathsf X_s=a+s,
\]
and scaling by a positive constant $k$ should be a map $\mathsf Y_k$ with
$a\circ\mathsf Y_k=ka$.  Two things are deliberately kept apart.

\begin{enumerate}
  \item The \emph{motions} record arithmetic propagation.  They are required to
  transport the assignment in a prescribed way, and they need not preserve the
  metric.
  \item The \emph{metric} measures the infinitesimal cost of moving.  It is part
  of the structure, and the axioms below constrain it through the assignment.
\end{enumerate}

The elementary model of Section~\ref{subsec:p0-basic-model} makes both visible at
once: its additive and multiplicative motions are affine maps of the upper
half-plane, and the metric is the invariant hyperbolic metric of the positive
affine group.

\subsection{Regular arithmetic expression spaces}

\begin{definition}[Regular arithmetic expression space]\label{def:regular-aes}
A \emph{regular arithmetic expression space}, or regular AES, is a tuple
\[
 \mathfrak E=(M,g,a;\mu,\lambda)
\]
such that $M$ is an oriented smooth surface, possibly with boundary, $g$ is a smooth Riemannian
metric, $a\in C^\infty(M,\mathbb R)$, $\mu\in\mathbb R\setminus\{0\}$,
$\lambda\in\mathbb R$, and
\begin{equation}\label{eq:regular-aes-eikonal}
 |\nabla a|_g^2=\mu^2+\lambda^2a^2
\end{equation}
everywhere on $M$.  The function $a$ is called the
\emph{assignment function}.
\end{definition}

The metric in Definition~\ref{def:regular-aes} is part of the structure.
Equation~\eqref{eq:regular-aes-eikonal} is therefore an eikonal equation on a
specified Riemannian surface, not a metric-free identity.  Its right-hand side
records the two arithmetic directions: the constant term $\mu^2$ is the
additive intensity and the assignment term $\lambda^2a^2$ the multiplicative
one.

Let $J$ denote rotation by $+\pi/2$ in the oriented tangent planes.  Since
$\mu\ne0$, the quantity
\[
 q=\sqrt{\mu^2+\lambda^2a^2}
\]
is everywhere positive.  The next proposition shows that the eikonal definition
contains exactly the oriented orthonormal directional data used in the flow
equation.

\begin{proposition}[Canonical arithmetic frame]\label{prop:canonical-arithmetic-frame}
For an oriented Riemannian surface $(M,g)$, a smooth function $a$, and
constants $\mu\ne0$ and $\lambda$, the following are equivalent.
\begin{enumerate}
\item The eikonal identity \eqref{eq:regular-aes-eikonal} holds.
\item There is a unique positively oriented orthonormal frame
      $(X_u,X_v)$ satisfying
      \begin{equation}\label{eq:arithmetic-frame-derivatives}
       X_u a=\mu,
       \qquad
       X_v a=\lambda a.
      \end{equation}
\end{enumerate}
When these conditions hold, put $E_a=\nabla a/q$.  The frame is
\begin{equation}\label{eq:explicit-arithmetic-frame}
 \boxed{
 \begin{aligned}
 X_u&=\frac{\mu}{q}E_a-\frac{\lambda a}{q}J E_a,\\
 X_v&=\frac{\lambda a}{q}E_a+\frac{\mu}{q}J E_a.
 \end{aligned}}
\end{equation}
\end{proposition}

\begin{proof}
Assume the eikonal identity.  Then $|\nabla a|=q>0$, so $E_a$ and $JE_a$ form a
smooth positively oriented orthonormal frame.  The coefficient matrix in
\eqref{eq:explicit-arithmetic-frame} is
\[
 \frac1q
 \begin{pmatrix}
  \mu&\lambda a\\
  -\lambda a&\mu
 \end{pmatrix},
\]
whose columns are orthonormal and whose determinant is one.  Hence
$(X_u,X_v)$ is positively oriented and orthonormal.  Since
$da(E_a)=q$ and $da(JE_a)=0$, the two identities in
\eqref{eq:arithmetic-frame-derivatives} follow.

For uniqueness, write any positively oriented orthonormal frame in the form
\[
 X_u=cE_a+sJE_a,
 \qquad
 X_v=-sE_a+cJE_a,
 \qquad c^2+s^2=1.
\]
The two directional identities force $qc=\mu$ and $-qs=\lambda a$.
Thus $c=\mu/q$ and $s=-\lambda a/q$, which gives
\eqref{eq:explicit-arithmetic-frame}.

Conversely, if an orthonormal frame satisfies
\eqref{eq:arithmetic-frame-derivatives}, then
\[
 |\nabla a|^2=(X_u a)^2+(X_v a)^2
 =\mu^2+\lambda^2a^2.
\]
\end{proof}

The frame is generally non-coordinate.  Applying its bracket to the assignment
function gives
\begin{equation}\label{eq:arithmetic-frame-bracket-on-a}
 [X_u,X_v]a
 =X_u(\lambda a)-X_v(\mu)
 =\mu\lambda.
\end{equation}
Thus, when $\lambda\ne0$, the fields cannot commute.  If
$(\eta^u,\eta^v)$ is their dual coframe, the valid intrinsic Pfaffian statement
is
\[
 da=\mu\,\eta^u+\lambda a\,\eta^v;
\]
the one-forms $\eta^u,\eta^v$ are not presumed exact.  Equation
\eqref{eq:arithmetic-frame-bracket-on-a} is the first appearance of a quantity
that returns in Sections~\ref{sec:p0-acs-tearing} and~\ref{sec:p0-contact}: the
additive direction changes the multiplicative datum by $\mu\lambda$ per unit
area of the arithmetic frame.

\subsection{The basic hyperbolic model}
\label{subsec:p0-basic-model}

Let
\[
 \mathbb H^2=\{(x,y)\in\mathbb R^2:y>0\}
\]
with its standard orientation, and fix $\mu,\lambda>0$.  Put
\begin{equation}\label{eq:p0-e0-data}
 g_{\mu,\lambda}
 =\frac1{y^2}
  \left(\frac{\dd x^2}{\mu^2}
       +\frac{\dd y^2}{\lambda^2}\right),
 \qquad
 a(x,y)=-\frac{x}{y},
\end{equation}
and record the assignment separately, since every later computation refers to
it:
\begin{equation}\label{eq:basic-hyperbolic-assignment}
 a(x,y)=-\frac{x}{y}.
\end{equation}
We call
\[
 \mathfrak E_0(\mu,\lambda)=(\mathbb H^2,g_{\mu,\lambda},a;\mu,\lambda)
\]
the \emph{basic hyperbolic model}.  In earlier drafts of the series and in
Paper~I this model was written $\mathfrak E_{\mathrm{hyp}}(\mu,\lambda)$; the
identification is stated once, here, and the letter $\mathfrak E_0$ is used
from now on (see also Subsection~\ref{subsec:p0-model-names}).

The metric admits the group-theoretic description that explains why addition
and multiplication are the natural motions.  Write
\[
 G=\operatorname{Aff}^{+}(1,\mathbb R)
   =\{(x,y):x\in\mathbb R,\ y>0\}
\]
with multiplication
\begin{equation}\label{eq:positive-affine-group-law}
 (x,y)(x',y')=(x+yx',yy').
\end{equation}
The element $(x,y)$ represents the affine map $t\mapsto x+yt$.  Left
translation by $(x_0,y_0)$ is
\[
 L_{(x_0,y_0)}(x,y)=(x_0+y_0x,y_0y).
\]
For fixed $\mu,\lambda>0$, consider the left-invariant vector fields
\begin{equation}\label{eq:hyperbolic-arithmetic-frame}
 X_u=-\mu y\,\partial_x,
 \qquad
 X_v=-\lambda y\,\partial_y .
\end{equation}
Both signs are chosen so that the frame is positively oriented for
$dx\wedge dy$ and has the action on the assignment stated below.

\begin{proposition}[Invariant affine metric]\label{prop:invariant-affine-metric}
The Riemannian metric for which the frame
\eqref{eq:hyperbolic-arithmetic-frame} is orthonormal is
\begin{equation}\label{eq:hyperbolic-aes-metric}
 g_{\mu,\lambda}
 =\frac1{y^2}
  \left(\frac{dx^2}{\mu^2}+\frac{dy^2}{\lambda^2}\right).
\end{equation}
It is invariant under every left translation of $G$.
\end{proposition}

\begin{proof}
The coframe dual to $(X_u,X_v)$ is
\[
 -\frac{dx}{\mu y},
 \qquad
 -\frac{dy}{\lambda y},
\]
whose sum of squares is \eqref{eq:hyperbolic-aes-metric}.  Under
$L_{(x_0,y_0)}$, one has
\[
 x\longmapsto x_0+y_0x,\qquad
 y\longmapsto y_0y,\qquad
 dx\longmapsto y_0dx,\qquad
 dy\longmapsto y_0dy.
\]
Substitution in \eqref{eq:hyperbolic-aes-metric} shows
$L_{(x_0,y_0)}^*g_{\mu,\lambda}=g_{\mu,\lambda}$.
\end{proof}

\begin{theorem}[Basic hyperbolic AES]\label{thm:basic-hyperbolic-aes}
For every $\mu,\lambda>0$, the tuple
\[
 \mathfrak E_0(\mu,\lambda)
 =(\mathbb H^2,g_{\mu,\lambda},a;\mu,\lambda)
\]
is a regular arithmetic expression space.  Its canonical arithmetic frame is
the left-invariant frame \eqref{eq:hyperbolic-arithmetic-frame}.  Moreover,
$(\mathbb H^2,g_{\mu,\lambda})$ is complete and has constant Gaussian
curvature
\[
 K=-\lambda^2.
\]
\end{theorem}

\begin{proof}
Direct differentiation gives
\[
 a_x=-\frac1y,
 \qquad
 a_y=\frac{x}{y^2}=-\frac ay.
\]
Therefore the frame \eqref{eq:hyperbolic-arithmetic-frame} satisfies
\[
 X_u a=\mu,
 \qquad
 X_v a=\lambda a.
\]
It is positively oriented and orthonormal by
Proposition~\ref{prop:invariant-affine-metric}, and the frame equivalence of
Proposition~\ref{prop:canonical-arithmetic-frame} proves the regular-AES
identity.  Equivalently, using
\[
 g_{\mu,\lambda}^{-1}
 =\mu^2y^2\,\partial_x\otimes\partial_x
  +\lambda^2y^2\,\partial_y\otimes\partial_y,
\]
one obtains directly
\[
 |\nabla a|_{g_{\mu,\lambda}}^2
 =\mu^2y^2a_x^2+\lambda^2y^2a_y^2
 =\mu^2+\lambda^2a^2.
\]

For completeness and curvature, introduce the global coordinate
\[
 X=\frac{\lambda}{\mu}x.
\]
Then
\begin{equation}\label{eq:scaled-standard-hyperbolic-metric}
 g_{\mu,\lambda}
 =\frac1{\lambda^2}\frac{dX^2+dy^2}{y^2}.
\end{equation}
The metric in parentheses is the complete standard upper-half-plane metric of
curvature $-1$.  Multiplication of a Riemannian metric by
$\lambda^{-2}$ preserves completeness and multiplies Gaussian curvature by
$\lambda^2$.  Hence $g_{\mu,\lambda}$ is complete and has curvature
$-\lambda^2$.
\end{proof}

The parameters have different geometric roles.  The change
$X=(\lambda/\mu)x$ removes $\mu$ from the metric, so $\mu$ fixes the
additive normalization of the assignment frame.  The parameter $\lambda$
sets the curvature scale.  No uniqueness claim for regular AESs follows from
this normalization.

\paragraph{Horocyclic coordinates and assignment contours.}
Set $y=e^{\lambda v}$ and $u=x$.  Then
\begin{equation}\label{eq:horocyclic-affine-coordinates}
 g_{\mu,\lambda}
 =e^{-2\lambda v}\frac{du^2}{\mu^2}+dv^2,
 \qquad
 a(u,v)=-u e^{-\lambda v}.
\end{equation}
The curves $y=\mathrm{constant}$ are horocycles based at the ideal point
$\infty$, and the vertical lines $x=\mathrm{constant}$ are their
orthogonal geodesics.  The assignment contour of value $c$ is
\[
 a^{-1}(c)=\{x=-cy,\ y>0\}.
\]
In particular, the zero locus in this model is the vertical geodesic
$\{x=0,\ y>0\}$.

\paragraph{The Laplace eigenfunction.}
We use the divergence-of-gradient sign convention
\begin{equation}\label{eq:laplacian-sign-convention}
 \Delta_g f
 =\frac1{\sqrt{\det g}}\,
  \partial_i\!\left(\sqrt{\det g}\,g^{ij}\partial_jf\right).
\end{equation}

\begin{proposition}[Laplace eigenfunction]\label{prop:laplace-eigenfunction}
The assignment function of the basic hyperbolic AES satisfies
\begin{equation}\label{eq:hyperbolic-assignment-laplacian}
 \Delta_{g_{\mu,\lambda}}a=2\lambda^2a.
\end{equation}
\end{proposition}

\begin{proof}
Since
\[
 \sqrt{\det g_{\mu,\lambda}}=\frac1{\mu\lambda y^2},
\]
the products
$\sqrt{\det g}\,g^{xx}=\mu/\lambda$ and
$\sqrt{\det g}\,g^{yy}=\lambda/\mu$ are constant.  Thus
\[
 \Delta_{g_{\mu,\lambda}}f
 =\mu^2y^2 f_{xx}+\lambda^2y^2 f_{yy}.
\]
For $a=-x/y$, one has $a_{xx}=0$ and
$a_{yy}=-2x/y^3=2a/y^2$.  Substitution gives
\eqref{eq:hyperbolic-assignment-laplacian}.
\end{proof}

\subsection{The four operations as motions}

On the basic model, define for $s\in\mathbb R$ and $k>0$
\begin{equation}\label{eq:p0-grid-moves}
  \mathsf X_s(x,y)=(x-sy,y),
  \qquad
  \mathsf Y_k(x,y)=\left(x,\frac yk\right).
\end{equation}

\begin{proposition}[Arithmetic action on the assignment]
\label{prop:p0-grid-assignment-action}
The maps in \eqref{eq:p0-grid-moves} satisfy
\[
  a\circ\mathsf X_s=a+s,
  \qquad
  a\circ\mathsf Y_k=ka.
\]
Consequently, translation and positive scaling of the current arithmetic value
can be represented by motions of points in $\mathfrak E_0$.
\end{proposition}

\begin{proof}
Direct substitution gives
\[
  a(x-sy,y)=-\frac{x-sy}{y}=a(x,y)+s,
\]
and
\[
  a\left(x,\frac yk\right)
  =-\frac{x}{y/k}=k a(x,y).
\]
\end{proof}

These transformations are assignment-compatible; they are not generally
isometries of $g_{\mu,\lambda}$.  The distinction is the one announced in
Subsection~\ref{subsec:p0-basic-model}: the motions record arithmetic
propagation, whereas the metric measures its infinitesimal geometry.

The remaining two operations are obtained from the same two families and a
reflection.

\begin{center}
\renewcommand{\arraystretch}{1.35}
\begin{tabular}{lll}
\toprule
operation & motion of $\mathfrak E_0$ & effect on $a$ \\
\midrule
addition of $s$ & $\mathsf X_s$ & $a\mapsto a+s$ \\
subtraction of $s$ & $\mathsf X_{-s}$ & $a\mapsto a-s$ \\
multiplication by $k>0$ & $\mathsf Y_k$ & $a\mapsto ka$ \\
division by $k>0$ & $\mathsf Y_{1/k}$ & $a\mapsto a/k$ \\
reflection (sign change) & $(x,y)\mapsto(-x,y)$ & $a\mapsto -a$ \\
\bottomrule
\end{tabular}
\end{center}

Subtraction is thus the inverse of the additive motion, and division by a
nonzero constant is the inverse of the multiplicative motion.  The reflection is
not by itself one of the four operations; it appears because the assignment of
the model is signed, and it is the elementary instance of the second-slot
subtraction $c-z$ of Section~\ref{sec:p0-expansion}.

For an integer $m\ge2$, the two moves satisfy the exact relation
\begin{equation}\label{eq:p0-bs-relation}
  \mathsf Y_m^{-1}\circ\mathsf X_s^m\circ\mathsf Y_m
  =\mathsf X_s.
\end{equation}
Indeed, $\mathsf X_s^m=\mathsf X_{ms}$, and conjugation by $\mathsf Y_m$
divides the additive displacement by $m$.  Equation
\eqref{eq:p0-bs-relation} is the elementary Baumslag--Solitar-type relation
visible in the grid \cite{BaumslagSolitar1962NonHopfian}; no complexity
conclusion is inferred from it.  Its arithmetic content is that a
multiplicative motion transports the additive ruler: after rescaling the
space, the additive unit that produces the same effect on $a$ is $1/m$ of the
original one.  The same transport, made infinitesimal and measured against a
commutative shadow, is the tearing discussed in
Section~\ref{sec:p0-acs-tearing}.

\paragraph{Division by the evolving value.}
The motions above are affine, and so is every word in them.  The remaining
elementary context of Section~\ref{sec:p0-expansion}, second-slot division
$c/z$, is not affine and cannot be realized by a motion of $\mathfrak E_0$ that
preserves the assignment line.  It is realized projectively: on the projective
line it is an element of $\PGL_2(K)$, and it moves the arithmetic value
$z\mapsto c/z$ on the ordinary domain $z\ne0$.  Section~\ref{sec:p0-ripples}
reads this operation not as a motion of the value but as a transport of
conditions, and Section~\ref{sec:p0-projective-unification} places both
readings in one matrix.

\subsection{Point paths}
\label{subsec:p0-point-paths}

The left-expanded comb becomes geometric once its one-hole sections are
restricted to translations and nonzero scalings.  Let a left-expanded affine
expression have chronological sections
\[
  F_i(z)=\alpha_i z+\beta_i,
  \qquad \alpha_i\ne0.
\]
For an initial arithmetic value $z_0$, define
\[
  z_i=F_i(z_{i-1}),
  \qquad i=1,\ldots,n.
\]
Choose a point $p_0\in\mathbb H^2$ with $a(p_0)=z_0$, and realize each
translation or positive scaling by the corresponding map in
\eqref{eq:p0-grid-moves}.  The successive points
\[
  p_0,p_1,\ldots,p_n
\]
form the \emph{point path} of the bounded expression in this chosen realization.
By Proposition~\ref{prop:p0-grid-assignment-action},
\[
  a(p_i)=z_i.
\]
Thus the geometry records every intermediate value, not only the endpoint.  The
word ``point path'' is a descriptive name for this orbit; it asserts no new
structure beyond the displayed construction.

\begin{example}[A four-step path]
\label{ex:p0-four-step-path}
Start from $z_0=1$ and apply
\[
  z\mapsto4z,
  \qquad z\mapsto z-1,
  \qquad z\mapsto2z,
  \qquad z\mapsto z-3.
\]
The value history is
\[
  1\longmapsto4\longmapsto3\longmapsto6\longmapsto3.
\]
In the grid this becomes two transverse scale moves and two horocyclic
translation moves.  The final value $3$ alone does not determine this path.
\end{example}

\begin{figure}[t]
\centering
\begin{tikzpicture}[
  x=0.62cm,y=0.58cm,
  axis/.style={-{Latex[length=2mm]},thick,black!70},
  hline/.style={thin,blue!45},
  vline/.style={thin,green!55!black},
  path/.style={very thick,black},
  every node/.style={font=\scriptsize}
]
  \fill[blue!2] (-8.3,0) rectangle (8.3,8.7);
  \draw[axis] (-8.5,0)--(8.7,0) node[right] {$x$};
  \draw[axis] (0,-0.1)--(0,9.0) node[above] {$y$};
  \foreach \yy in {1,2,4,8}
    {\draw[hline] (-8.2,\yy)--(8.2,\yy);}
  \foreach \xx in {-8,-6,-4,-3,-2,-1,1,2,3,4,6,8}
    {\draw[vline] (\xx,0.08)--(\xx,8.45);}
  \node[red!65!black] at (-8.0,8.35) {$1$};
  \node[red!65!black] at (-8.0,4.25) {$2$};
  \node[red!65!black] at (-8.0,2.25) {$4$};
  \node[red!65!black] at (-8.0,1.25) {$8$};
  \node[red!65!black] at (0.25,8.35) {$0$};

  \coordinate (p0) at (-8,8);
  \coordinate (p1) at (-8,2);
  \coordinate (p2) at (-6,2);
  \coordinate (p3) at (-6,1);
  \coordinate (p4) at (-3,1);
  \draw[path] (p0)--(p1)--(p2)--(p3)--(p4);
  \foreach \p/\lab/\pos in {p0/{1}/above,p1/{4}/left,p2/{3}/above,p3/{6}/left,p4/{3}/above}
    {\fill (\p) circle (2pt) node[\pos] {$\lab$};}
  \node[font=\small,anchor=west] at (-1.6,7.3)
    {$1\xrightarrow{\times4}4\xrightarrow{-1}3
      \xrightarrow{\times2}6\xrightarrow{-3}3$};
\end{tikzpicture}
\caption{A schematic point path in the basic hyperbolic arithmetic grid.  The
black polyline retains the intermediate propagation that the endpoint value
forgets.}
\label{fig:p0-path-grid}
\end{figure}

\paragraph{The affine algebra behind the path.}
Every affine prefix has a unique form
\[
  G_i(z)=A_i z+B_i,
  \qquad A_i\ne0.
\]
If $F_i(z)=\alpha_i z+\beta_i$, then
\begin{equation}\label{eq:p0-affine-prefix-recursion}
  A_i=\alpha_i A_{i-1},
  \qquad
  B_i=\alpha_i B_{i-1}+\beta_i,
  \qquad (A_0,B_0)=(1,0).
\end{equation}
The point path is therefore the evaluation of an affine matrix path
\[
  \begin{pmatrix}A_i&B_i\\0&1\end{pmatrix}
  \begin{pmatrix}z_0\\1\end{pmatrix}.
\]
The lower-left entry remains zero at every stage.  Projectively this means that
all prefixes stabilize the same point $\infty$.  Section~\ref{sec:p0-ripples}
shows what changes when right-slot division makes that entry nonzero.

\subsection{Continuous arithmetic motion}
\label{subsec:p0-continuous-motion}

The discrete motions of Subsection~\ref{subsec:p0-point-paths} have an
infinitesimal counterpart.  Translation and dilation generate the Lie algebra
of the positive affine group, and their linear combinations give the continuous
propagation law.  The formulation below is intrinsic: the symbols $u$ and $v$
label arithmetic directions, but need not be coordinates.

\paragraph{The affine Lie algebra and its flow.}
On the assignment line with coordinate $z$, let
\[
 E=\partial_z,
 \qquad
 H=z\partial_z .
\]
These are respectively the infinitesimal generators of translation and positive
dilation, and they satisfy $[E,H]=E$.  Fix constants
$\mu,\lambda\in\mathbb R$, with $\mu\ne0$.  The generator selected by an angle
$\theta$ is
\begin{equation}\label{eq:angular-affine-generator}
 \Omega_\theta
 =\mu\cos\theta\,E+\lambda\sin\theta\,H .
\end{equation}
Consequently, an integral curve $z=z(s)$ satisfies
\begin{equation}\label{eq:continuous-affine-flow}
 \frac{dz}{ds}
 =\mu\cos\theta+\lambda z\sin\theta .
\end{equation}
This Lie-algebra calculation is a derivation of the continuous form of the
arithmetic motion, and it fixes the convention: the translation coefficient is
$\mu\cos\theta$, while the dilation coefficient is $\lambda\sin\theta$.

When $\theta$ is constant, set $b=\lambda\sin\theta$.  The solution through
$z(0)=z_0$ is
\begin{equation}\label{eq:directional-affine-solution}
 z(s)=
 \begin{cases}
  e^{bs}z_0+\displaystyle\mu\cos\theta\,\frac{e^{bs}-1}{b},
     & b\ne0,\\[7pt]
  z_0+\mu s\cos\theta, & b=0.
 \end{cases}
\end{equation}
In particular, the positive arithmetic axes give
\[
 \theta=0:\quad z(s)=z_0+\mu s,
 \qquad
 \theta=\frac{\pi}{2}:\quad z(s)=e^{\lambda s}z_0.
\]
The opposite axes give subtraction and inverse dilation.

\paragraph{The Pfaffian propagation constraint.}
There is a useful three-dimensional local state space with coordinates
$(u,v,z)$.  On it define
\begin{equation}\label{eq:affine-pfaffian-form}
 \alpha=dz-\mu\,du-\lambda z\,dv
\end{equation}
and the two horizontal vector fields
\begin{equation}\label{eq:affine-horizontal-fields}
 D_u=\partial_u+\mu\partial_z,
 \qquad
 D_v=\partial_v+\lambda z\partial_z.
\end{equation}
A curve tangent to
$\cos\theta\,D_u+\sin\theta\,D_v$ satisfies
\eqref{eq:continuous-affine-flow}.  Thus the frequently used notation
\[
 dz=\mu\,du+\lambda z\,dv
\]
means that $\alpha$ vanishes on admissible tangent vectors.  It is not an
identity asserting that $z$ is a function of two independent coordinates whose
ordinary partial derivatives are simultaneously $\mu$ and $\lambda z$.
Indeed,
\begin{equation}\label{eq:horizontal-noncommutation}
 [D_u,D_v]=\mu\lambda\,\partial_z.
\end{equation}
When $\mu\lambda\ne0$, this bracket already rules out such an integrable
coordinate interpretation.  The same bracket is computed intrinsically in
\eqref{eq:arithmetic-frame-bracket-on-a}; the local model here and the frame
there describe one object.  Its geometric reading as a contact structure is
taken up in Section~\ref{sec:p0-contact}.

\begin{theorem}[Continuous affine flow]\label{thm:continuous-affine-flow}
Let $\mathfrak E=(M,g,a;\mu,\lambda)$ be a regular AES with its canonical
arithmetic frame.  If a unit-speed curve $\gamma$ has tangent
\[
 \dot\gamma
 =\cos\theta\,X_u+\sin\theta\,X_v,
\]
where $\theta$ may vary along the curve, then
\begin{equation}\label{eq:intrinsic-continuous-affine-flow}
 \frac{d}{ds}(a\circ\gamma)
 =\mu\cos\theta+\lambda(a\circ\gamma)\sin\theta.
\end{equation}
For constant $\theta$, its solutions are given by \eqref{eq:directional-affine-solution}.
\end{theorem}

\begin{proof}
By the chain rule and \eqref{eq:arithmetic-frame-derivatives},
\[
 \frac{d}{ds}(a\circ\gamma)
 =da(\dot\gamma)
 =\cos\theta\,X_u a+\sin\theta\,X_v a
 =\mu\cos\theta+\lambda a\sin\theta.
\]
For constant $\theta$, this is the linear equation already integrated in
\eqref{eq:directional-affine-solution}.
\end{proof}

\paragraph{Rectification.}
For $\lambda\ne0$, define
\begin{equation}\label{eq:rectifying-assignment}
 r=\operatorname{arcsinh}\!\left(\frac{\lambda a}{\mu}\right).
\end{equation}
The chain rule and \eqref{eq:regular-aes-eikonal} give
\begin{equation}\label{eq:rectified-eikonal}
 |\nabla r|=|\lambda|.
\end{equation}
Thus $r$ has constant gradient magnitude.  When $\lambda=0$, the original
equation already reduces to $|\nabla a|=|\mu|$, and $a/\mu$ is a
unit-gradient rectifying function.  Rectification of the assignment is
possible; rectification of *both* arithmetic directions at once is not, by
\eqref{eq:horizontal-noncommutation}.  This asymmetry is the discrete root of
the tearing language of Section~\ref{sec:p0-holed}.

The same affine convention fixes the sign of the elementary local order defect.
For small increments $h$ and $k$, let
\[
 A_h(z)=z+\mu h,
 \qquad
 M_k(z)=e^{\lambda k}z.
\]
Then
\begin{align}
 (M_k\circ A_h)(z)-(A_h\circ M_k)(z)
 &=\mu h\bigl(e^{\lambda k}-1\bigr)\notag\\
 &=\mu\lambda\,hk+O\!\left(\lVert(h,k)\rVert^3\right).
 \label{eq:infinitesimal-order-defect}
\end{align}
This is an infinitesimal two-order comparison, not an exact area formula for an
arbitrary finite region.  The exact finite statement, and its interpretation as
a defect measured in the commutativized space, occupy
Section~\ref{sec:p0-acs-tearing}.

\paragraph{Affine motion inside projective motion.}
An infinitesimal projective transformation of the line has the Riccati form
\[
 \dot z=\beta+\alpha z+\kappa z^2.
\]
The theory developed in this section is the affine slice $\kappa=0$, with
$\beta=\mu\cos\theta$ and $\alpha=\lambda\sin\theta$.  Results proved for
regular affine AESs do not automatically extend to the quadratic projective
flow.

\subsection{Model names: $E_0$, $E_1$, and later examples}
\label{subsec:p0-model-names}

The historical drafts of this programme used the labels $E_0$ and $E_1$
inconsistently; in some notes the hyperbolic half-plane is called $E_0$, in
others $E_1$.  Because later chapters introduce spaces that are *not* smooth,
we fix the convention now, once, and record it as a convention rather than a
theorem.

\begin{convention}[Model labels]
\label{conv:p0-model-labels}
$\mathfrak E_0(\mu,\lambda)$ denotes the basic regular hyperbolic model of
Subsection~\ref{subsec:p0-basic-model}.  For $k\ge1$, the label
$\mathfrak E_k$ is used descriptively for an arithmetic expression space whose
regular locus is a surface with $k$ declared punctures: points removed from the
regular structure and carrying no regular AES data.  No classification by $k$
is asserted; the label records the number of declared punctures and nothing
else.  The reserved historical labels $E_k$ and $E_{\log}$ of earlier drafts
are not model classes, and this paper defines neither of them.
\end{convention}

Two examples make the convention concrete; both are developed later, and the
first one is the only example in this paper whose punctured model is verified
in detail.

\begin{example}[The basic model, $k=0$]
\label{ex:p0-e0-model}
$\mathfrak E_0(\mu,\lambda)$ is the upper half-plane with the metric
\eqref{eq:p0-e0-data} and the assignment $a=-x/y$.  It has no puncture, its
assignment vanishes along one complete geodesic, and by
Theorem~\ref{thm:basic-hyperbolic-aes} it is complete with constant curvature.
This is the model in which the arithmetic motions of
Subsection~\ref{subsec:p0-point-paths} are read as point paths.
\end{example}

\begin{example}[A one-puncture model, $k=1$]
\label{ex:p0-e1-model}
Let $\mathbb D=\{\rho<1\}$ be the unit disc with the complete Poincar\'e
metric $g_{\mathbb D}=4\,(d\xi^2+d\eta^2)/\lambda^2(1-\rho^2)^2$ and the radial
assignment
\[
 a_{\mathbb D}(\xi,\eta)=
 \begin{cases}
  \dfrac{2\mu}{\lambda}\,\dfrac{\rho}{1-\rho^2},&\rho\ne0,\\[6pt]
  0,&\rho=0 .
 \end{cases}
\]
On the punctured disc $\mathbb D\setminus\{0\}$ this is a regular AES with
$|\nabla a_{\mathbb D}|^2=\mu^2+\lambda^2a_{\mathbb D}^2$, while at the centre
the assignment fails to be $C^1$ although the metric extends smoothly.  The
regular locus is therefore $\mathbb D\setminus\{0\}$ and the model carries
exactly one puncture.  Paper~I verifies the details of this model as its
isolated-zero example; the global invariant $\lambda a_{\mathbb D}
=\mu\sinh(\lambda s)$ along hyperbolic distance $s$ from the centre shows that
the assignment, not the metric, is what fails at the puncture.
\end{example}

Example~\ref{ex:p0-e1-model} is the first holed arithmetic expression space, and
Section~\ref{sec:p0-holed} returns to it.  Its puncture is not a hole in the
metric; it is a point at which the regular structure stops being defined, and
it is where the arithmetic motions of the punctured part can no longer be
extended.
